\begin{frame}
  \frametitle{Diamond Differencing}
  To close \cref{eq:fv-eqs}, we need an expression for the cell fluxes.
  Recall the forward and backward finite difference approximations:
  \begin{subequations}
    \begin{equation}
      \label{eq:forward-diff}
      \dv{\psi}{x}\left. \right|_{i,j} = \frac{\psi^{i+1/2}-\psi^i}{h}+\mathcal{O}\left(h\right)
    \end{equation}
    \begin{equation}
      \label{eq:backward-diff}
      \dv{\psi}{x}\left. \right|_{i,j} = \frac{\psi^{i}-\psi^{i-1/2}}{h}+\mathcal{O}\left(h\right)
    \end{equation}
  \end{subequations}
  Subtracting \cref{eq:forward-diff} and \cref{eq:backward-diff} gives:
  \begin{equation}
    0 = \frac{\psi^{i+1/2}-2\psi^i+\psi^{i-1/2}}{h}+\mathcal{O}\left(h\right)
  \end{equation}
  Or, equivalently:
  \begin{equation}\label{eq:diamond-differencing}
    \psi^i = \frac{\psi^{i+1/2}+\psi^{i-1/2}}{2}+\mathcal{O}\left(h^2\right)
  \end{equation}
  \cref{eq:diamond-differencing} is known as the \textit{diamond differencing} closure and is 2\textsuperscript{nd}-order accurate in space.
  Depending on the available information from surrounding cells (i.e. the sign of $\mu$ and $\eta$), \cref{eq:diamond-differencing} is used to create a \textit{sweeping} algorithm that solves \cref{eq:fv-eqs} along the direction of travel.
\end{frame}